% from ega3-1-fr.tex
\oldpage[III]{166}

\S~13. Limites projectives en algèbre homologique \dotfill 64

13.1. La condition de Mittag-Leffler \dotfill 64

13.2. La condition de Mittag-Leffler pour les groupes abéliens \dotfill 65

13.3. Application~: cohomologie d'une limite projective de faisceaux \dotfill 68

13.4. Condition de Mittag-Leffler et objets gradués associés aux systèmes
projectifs \dotfill 69

13.5. Limites projectives de suites spectrales de complexes filtrés \dotfill 71

13.6. Suite spectrale d'un foncteur relative à un objet muni d'une filtration
finie \dotfill 73

13.7. Foncteurs dérivés d'une limite projective d'arguments \dotfill 75

\medskip
\textsc{Chapitre III.} --- \textbf{Étude cohomologique des faisceaux
cohérents} \dotfill 81

\medskip
\S~1. Cohomologie des schémas affines \dotfill 82

1.1. Rappels sur le complexe de l'algèbre extérieure \dotfill 82

1.2. Cohomologie de Čech d'un recouvrement ouvert \dotfill 85

1.3. Cohomologie d'un schéma affine \dotfill 88

1.4. Application à la cohomologie des préschémas quelconques \dotfill 89

\medskip
\S~2. Étude cohomologique des morphismes projectifs \dotfill 95

2.1. Calculs explicites de certains groupes de cohomologie \dotfill 95

2.2. Le théorème fondamental des morphismes projectifs \dotfill 100

2.3. Application aux faisceaux gradués d'algèbres et de modules \dotfill 102

2.4. Une généralisation du théorème fondamental \dotfill 107

2.5. Caractéristique d'Euler-Poincaré et polynôme de Hilbert \dotfill 109

2.6. Application~: critères d'amplitude \dotfill 111

\medskip
\S~3. Le théorème de finitude pour les morphismes propres \dotfill 115

3.1. Le lemme de dévissage \dotfill 115

3.2. Le théorème de finitude~: cas des schémas usuels \dotfill 116

3.3. Généralisation du théorème de finitude (schémas usuels) \dotfill 118

3.4. Le théorème de finitude~: cas des schémas formels \dotfill 119

\medskip
\S~4. Le théorème fondamental des morphismes propres. Applications \dotfill 122

4.1. Le théorème fondamental \dotfill 122

4.2. Cas particuliers et variantes \dotfill 129

4.3. Le théorème de connexion de Zariski \dotfill 130

4.4. Le «~main theorem~» de Zariski \dotfill 135

4.5. Complétés de modules d'homomorphismes \dotfill 138

4.6. Relations entre morphismes formels et morphismes usuels \dotfill 139

4.7. Un critère d'amplitude \dotfill 145

4.8. Morphismes finis de préschémas formels \dotfill 146

